\Chapter{41}

\Verse{1}
{
    \zA{话说刘老老两只手比着说道：“花儿落了结个大倭瓜。”}
}
{
    \eA{Old goody Liu, so the story goes, exclaimed, while making signs with both
        hands, "The flower dropped and a huge melon formed;" to the intense
        amusement of all the inmates,}
}
{
}

\Verse{2}
{
    \zA{众人听了，哄堂大笑 起来。}
}
{
    \eA{who burst into a boisterous fit of laughter. In due course, however, she
        drank the closing cup.}
}
{
}

\Verse{3}
{
    \zA{于是吃过门杯，因又斗趣笑道：“今儿实说罢，我的手脚子粗，又喝了酒， 仔细失手打了这磁杯。}
}
{
    \eA{Then she made another effort to evoke merriment. "To speak the truth
        to-day," she smilingly observed, "my hands and my feet are so rough, and
        I've had so much wine that I must be careful; or else I might, by a slip of
        the hand, break the porcelain cups.}
}
{
}

\Verse{4}
{
    \zA{有木头的杯取个来，我就失了手，掉了地下也无碍。”}
}
{
    \eA{If you have got any wooden cups, you'd better produce them. It wouldn't
        matter then if even they were to slip out of my hands and drop on the
        ground!" This joke excited some more mirth.}
}
{
}

\Verse{5}
{
    \zA{众人 听了又笑起来。}
}
{
    \eA{But lady Feng, upon hearing this speedily put on a smile.}
}
{
}

\Verse{6}
{
    \zA{凤姐儿听如此说，便忙笑道：“果真要木头的，我就取了来，可有 一句话先说下：这木头的可比不得磁的，那都是一套，定要吃遍一套才算呢。”}
}
{
    \eA{"Well," she said, "if you really want a wooden one, I'll fetch you one at
        once! But there's just one word I'd like to tell you beforehand. Wooden cups
        are not like porcelain ones. They go in sets; so you'll have to do the right
        thing and drink from every cup of the set." "I just now simply spoke in jest
        about those cups in order to induce them to laugh,"}
}
{
}

\Verse{7}
{
    \zA{刘 老老听了，心下？？}
}
{
    \eA{old goody Liu at these words, mused within herself, "but,}
}
{
}

\Verse{8}
{
    \zA{道：“我方才不过是趣话取笑儿，谁知他果真竟有。}
}
{
    \eA{who would have thought that she actually has some of the kind. I've often
        been to the large households of village gentry on a visit,}
}
{
}

\Verse{9}
{
    \zA{我时常在 乡绅大家也赴过席，金杯银杯倒都也见过，从没见有木头杯的。}
}
{
    \eA{and even been to banquets there and seen both gold cups and silver cups; but
        never have I beheld any wooden ones about! Ah, of course! They must, I
        expect, be the wooden bowls used by the young children.}
}
{
}

\Verse{10}
{
    \zA{哦是了，想必是小 孩子们使的木碗儿，不过诓我多喝两碗。}
}
{
    \eA{Their object must be to inveigle me to have a couple of bowlfuls more than
        is good for me! But I don't mind it. This wine is, verily, like honey, so if
        I drink a little more,}
}
{
}

\Verse{11}
{
    \zA{别管他，横竖这酒蜜水儿似的，多喝点子 也无妨。”}
}
{
    \eA{it won't do me any harm." Bringing this train of thought to a close, "Fetch
        them!" she said aloud. "We'll talk about them by and bye."}
}
{
}

\Verse{12}
{
    \zA{想毕，便说“取来再商量”。}
}
{
    \eA{Lady Feng then directed Feng Erh to go and bring the set of ten cups,}
}
{
}

\Verse{13}
{
    \zA{凤姐因命丰儿：“前面里间书架子上，有 十个竹根套杯取来。”}
}
{
    \eA{made of bamboo roots, from the book-case in the front inner room. Upon
        hearing her orders, Feng Erh was about to go and execute them, when Yüan
        Yang smilingly interposed.}
}
{
}

\Verse{14}
{
    \zA{丰儿听了才要去取，鸳鸯笑道：“我知道，你那十个杯还小； 况且你才说木头的，这会子又拿了竹根的来，倒不好看。}
}
{
    \eA{"I know those ten cups of yours," she remarked, "they're small. What's more,
        a while back you mentioned wooden ones, and if you have bamboo ones brought
        now, it won't look well; so we'd better get from our place that set of ten
        large cups, scooped out of whole blocks of aspen roots,}
}
{
}

\Verse{15}
{
    \zA{不如把我们那里的黄杨根 子整？}
}
{
    \eA{and pour the contents of all ten of them down her throat?" "Yes, that would
        be much better,"}
}
{
}

\Verse{16}
{
    \zA{的十个大套杯拿来，灌他十下子。”}
}
{
    \eA{lady Feng smiled. The cups were then actually brought by a servant, at the
        direction of Yüan Yang.}
}
{
}

\Verse{17}
{
    \zA{凤姐儿笑道：“更好了。”}
}
{
    \eA{At the sight of them, old goody Liu was filled with surprise as well as with
        admiration.}
}
{
}

\Verse{18}
{
    \zA{鸳鸯果命人取来。}
}
{
    \eA{Surprise, as the ten formed one set going in gradation from large to small;}
}
{
}

\Verse{19}
{
    \zA{刘老老一看，又惊又喜：惊的是一连十个挨次大小分下来， 那大的足足的像个小盆子，极小的还有手里的杯子两个大；喜的是雕镂奇绝，一色
        山水树木人物，并有草字以及图印。}
}
{
    \eA{the largest being amply of the size of a small basin, the smallest even
        measuring two of those she held in her hand. Admiration, as they were all
        alike, engraved, in perfect style, with scenery, trees, and human beings,
        and bore inscriptions in the 'grass' character as well as the seal of the
        writer.}
}
{
}

\Verse{20}
{
    \zA{因忙说道：“拿了那小的来就是了。”}
}
{
    \eA{"It will be enough," she consequently shouted with alacrity, "if you give me
        that small one."}
}
{
}

\Verse{21}
{
    \zA{凤姐儿 笑道：“这个杯，没有这大量的，所以没人敢使他。}
}
{
    \eA{"There's no one," lady Feng laughingly insinuated, "with the capacity to
        tackle these! Hence it is that not a soul can pluck up courage enough to use
        them!}
}
{
}

\Verse{22}
{
    \zA{老老既要，好容易找出来，必 定要挨次吃一遍才使得。”}
}
{
    \eA{But as you, old dame, asked for them, and they were fished out, after ever
        so much trouble, you're bound to do the proper thing and drink out of each,}
}
{
}

\Verse{23}
{
    \zA{刘老老吓的忙道：“这个不敢!好姑奶奶，饶了我罢。”}
}
{
    \eA{one after the other." Old goody Liu was quite taken aback. "I daren't!" she
        promptly demurred. "My dear lady, do let me off!"}
}
{
}

\Verse{24}
{
    \zA{贾母、薛姨妈、王夫人知道他有年纪的人，禁不起，忙笑道：“说是说，笑是笑， 不可多吃了，只吃这头一杯罢。”}
}
{
    \eA{Dowager lady Chia, Mrs. Hsüeh and Madame Wang were quite alive to the fact
        that a person advanced in years as she was could not be gifted with such
        powers of endurance, and they hastened to smilingly expostulate. "To speak
        is to speak, and a joke is a joke, but she mayn't take too much," they said;
        "let her just empty this first cup, and have done." "O-mi-to-fu!"}
}
{
}

\Verse{25}
{
    \zA{刘老老道：“阿弥陀佛!我还是小杯吃罢，把这 大杯收着，我带了家去，慢慢的吃罢。”}
}
{
    \eA{ejaculated old goody Liu. "I'll only have a small cupful, and put this huge
        fellow away, and take it home and drink at my leisure." At this remark, the
        whole company once more gave way to laughter. Yüan Yang had no alternative
        but to give in and she had to bid a servant fill a large cup full of wine.}
}
{
}

\Verse{26}
{
    \zA{说的众人又笑起来。}
}
{
    \eA{Old goody Liu laid hold of it with both hands and raised it to her mouth.}
}
{
}

\Verse{27}
{
    \zA{鸳鸯无法，只得命人满斟了一大杯，刘老老两手捧着喝。}
}
{
    \eA{"Gently a bit!" old lady Chia and Mrs. Hsüeh shouted. "Mind you don't
        choke!" Mrs. Hsüeh then told lady Feng to put some viands before her. "Goody
        Liu!" smiled lady Feng,}
}
{
}

\Verse{28}
{
    \zA{贾母薛姨妈都道：“慢 些，别呛了。”}
}
{
    \eA{"tell me the name of anything you fancy, and I'll bring it and feed you."
        "What names can I know?"}
}
{
}

\Verse{29}
{
    \zA{薛姨妈又命凤姐儿布个菜儿。}
}
{
    \eA{old goody Liu rejoined. "Everything is good!" "Bring some egg-plant and
        salt-fish for her!"}
}
{
}

\Verse{30}
{
    \zA{凤姐笑道：“老老要吃什么，说出名 儿来，我夹了喂你。”}
}
{
    \eA{dowager lady Chia suggested with a smile. Lady Feng, upon hearing this
        suggestion, complied with it by catching some egg-plant and salt-fish with
        two chopsticks and putting them into old goody Liu's mouth.}
}
{
}

\Verse{31}
{
    \zA{刘老老道：“我知道什么名儿!样样都是好的。”}
}
{
    \eA{"You people," she smiled, "daily feed on egg-plants; so taste these of ours
        and see whether they've been nicely prepared or not."}
}
{
}

\Verse{32}
{
    \zA{贾母笑道： “把茄鲞夹些喂他。”}
}
{
    \eA{"Don't be making a fool of me!" old goody Liu answered smilingly. "If
        egg-plants can have such flavour,}
}
{
}

\Verse{33}
{
    \zA{凤姐儿听说，依言夹些茄鲞送入刘老老口中，因笑道：“你 们天天吃茄子，也尝尝我们这茄子，弄的可口不可口。”}
}
{
    \eA{we ourselves needn't sow any cereals, but confine ourselves to growing
        nothing but egg-plants!" "They're really egg-plants!" one and all protested.
        "She's not pulling your leg!" Old goody Liu was amazed. "If these be
        actually egg-plants," she said, "I've uselessly eaten them so long!}
}
{
}

\Verse{34}
{
    \zA{刘老老笑道：“别哄我了， 茄子跑出这个味儿来了，我们也不用种粮食，只种茄子了。”}
}
{
    \eA{But, my lady, do give me a few more; I'd like to taste the next mouthful
        carefully!" Lady Feng brought her, in very deed, another lot, and put it in
        her mouth. Old goody Liu munched for long with particular care.}
}
{
}

\Verse{35}
{
    \zA{众人笑道：“真是茄 子，我们再不哄你。”}
}
{
    \eA{"There is, it's true, something about them of the flavour of egg-plant," she
        laughingly remarked,}
}
{
}

\Verse{36}
{
    \zA{刘老老诧异道：“真是茄子？}
}
{
    \eA{"yet they don't quite taste like egg-plants. But tell me how they're cooked,}
}
{
}

\Verse{37}
{
    \zA{我白吃了半日。}
}
{
    \eA{so that I may prepare them in the same way for myself."}
}
{
}

\Verse{38}
{
    \zA{姑奶奶再喂我 些，这一口细嚼嚼。”}
}
{
    \eA{"There's nothing hard about it!" lady Feng answered smiling. "You take the
        newly cut egg-plants and pare the skin off.}
}
{
}

\Verse{39}
{
    \zA{凤姐儿果又夹了些放入他口内。}
}
{
    \eA{All you want then is some fresh meat. You hash it into fine mince, and fry
        it in chicken fat.}
}
{
}

\Verse{40}
{
    \zA{刘老老细嚼了半日，笑道： “虽有一点茄子香，只是还不像是茄子。}
}
{
    \eA{Then you take some dry chicken meat, and mix it with mushrooms, new bamboo
        shoots, sweet mushrooms, dry beancurd paste, flavoured with five spices,}
}
{
}

\Verse{41}
{
    \zA{告诉我是个什么法子弄的，我也弄着吃 去。”}
}
{
    \eA{and every kind of dry fruits, and you chop the whole lot into fine pieces.
        You then bake all these things in chicken broth,}
}
{
}

\Verse{42}
{
    \zA{凤姐儿笑道：“这也不难。}
}
{
    \eA{until it's absorbed, when you fry them, to finish, in sweet oil,}
}
{
}

\Verse{43}
{
    \zA{你把才下来的茄子把皮刨了，只要净肉，切成碎 钉子，用鸡油炸了。}
}
{
    \eA{and adding some oil, made of the grains of wine, you place them in a
        porcelain jar, and close it hermetically. At any time that you want any to
        eat, all you have to do is to take out some,}
}
{
}

\Verse{44}
{
    \zA{再用鸡肉脯子合香菌、新笋、蘑菇、五香豆腐干子、各色干果 子，都切成钉儿，拿鸡汤煨干了，拿香油一收，外加糟油一拌，盛在磁罐子里封严 了。}
}
{
    \eA{and mix it with some roasted chicken, and there it is all ready." Old goody
        Liu a shook her head and put out her tongue. "My Buddha's ancestor!" she
        shouted. "One wants about ten chickens to prepare this dish! It isn't
        strange then that it has this flavour!"}
}
{
}

\Verse{45}
{
    \zA{要吃的时候儿，拿出来，用炒的鸡瓜子一拌，就是了。”}
}
{
    \eA{Saying this, she quietly finished her wine. But still she kept on minutely
        scrutinizing the cup. "Haven't you yet had enough to satisfy you?"}
}
{
}

\Verse{46}
{
    \zA{刘老老听了，摇头吐 舌说：“我的佛祖!倒得多少只鸡配他，怪道这个味儿。”}
}
{
    \eA{lady Feng smiled. "If you haven't, well, then drink another cup."
        "Dreadful!" eagerly exclaimed old goody Liu. "I shall be soon getting so
        drunk that it will be the very death of me.}
}
{
}

\Verse{47}
{
    \zA{一面笑，一面慢慢的吃 完了酒，还只管细玩那杯子。}
}
{
    \eA{I was only looking at it as I admire pretty things like this! But what a
        trouble it must have cost to turn out!" "Have you done with your wine?"}
}
{
}

\Verse{48}
{
    \zA{凤姐笑道：“还不足兴，再吃一杯罢？”}
}
{
    \eA{Yuan Yang laughingly inquired. "But, after all, what kind of wood is this
        cup made of?"}
}
{
}

\Verse{49}
{
    \zA{刘老老忙道： “了不得，那就醉死了。}
}
{
    \eA{"It isn't to be wondered at," old goody Liu smiled, "that you can't make it
        out Miss! How ever could you people,}
}
{
}

\Verse{50}
{
    \zA{我因为爱这样儿好看，亏他怎么做来着！”}
}
{
    \eA{who live inside golden doors and embroidered apartments, know anything of
        wood! We have the whole day long the trees in the woods as our neighbours.}
}
{
}

\Verse{51}
{
    \zA{鸳鸯笑道：“酒 喝完了，到底这杯子是什么木头的？”}
}
{
    \eA{When weary, we use them as our pillows and go to sleep on them. When
        exhausted, we sit with our backs leaning against them. When, in years of
        dearth,}
}
{
}

\Verse{52}
{
    \zA{刘老老笑道：“怨不得姑娘不认得，你们在 这金门绣户里，那里认的木头？}
}
{
    \eA{we feel the pangs of hunger, we also feed on them. Day after day, we see
        them with our eyes; day after day we listen to them with our ears; day after
        day, we talk of them with our mouths.}
}
{
}

\Verse{53}
{
    \zA{我们成日家和树林子做街坊，困了枕着他睡，乏了 靠着他坐，荒年间饿了还吃他；眼睛里天天见他，耳朵里天天听他，嘴儿里天天说 他，所以好歹真假，我是认得的。}
}
{
    \eA{I am therefore well able to tell whether any wood be good or bad, genuine or
        false. Do let me then see what it is!" As she spoke, she intently scanned
        the cup for a considerable length of time. "Such a family as yours," she
        then said, "could on no account own mean things! Any wood that is easily
        procured,}
}
{
}

\Verse{54}
{
    \zA{让我认认。”}
}
{
    \eA{wouldn't even find a place in here.}
}
{
}

\Verse{55}
{
    \zA{一面说，一面细细端详了半日，道： “你们这样人家，断没有那贱东西，那容易得的木头你们也不收着了。}
}
{
    \eA{This feels so heavy, as I weigh it in my hands, that if it isn't aspen, it
        must, for a certainty, be yellow cedar." Her rejoinder amused every one in
        the room. But they then perceived an old matron come up. After asking
        permission of dowager lady Chia to speak: "The young ladies," she said,}
}
{
}

\Verse{56}
{
    \zA{我掂着这么 体？}
}
{
    \eA{"have got to the Lotus Fragrance pavilion,}
}
{
}

\Verse{57}
{
    \zA{，这再不是杨木，一定是黄松做的。”}
}
{
    \eA{and they request your commands, as to whether they should start with the
        rehearsal at once or tarry a while." "I forgot all about them!"}
}
{
}

\Verse{58}
{
    \zA{众人听了，哄堂大笑起来。}
}
{
    \eA{old lady Chia promptly cried with a smile. "Tell them to begin rehearsing at
        once!"}
}
{
}

\Verse{59}
{
    \zA{只见一个婆子走来，请问贾母说：“姑娘们都到了藕香榭，请示下：就演罢， 还是再等一会儿呢？”}
}
{
    \eA{The matron expressed her obedience and walked away. Presently, became
        audible the notes of the pan-pipe and double flute, now soft, now loud, and
        the blended accents of the pipe and fife. So balmy did the breeze happen to
        be and the weather so fine that the strains of music came wafted across the
        arbours and over the stream, and, needless to say, conduced to exhilarate
        their spirits and to cheer their hearts.}
}
{
}

\Verse{60}
{
    \zA{贾母忙笑道：“可是倒忘了，就叫他们演罢。”}
}
{
    \eA{Unable to resist the temptation, Pao-yü was the first to snatch a decanter
        and to fill a cup for himself. He quaffed it with one breath.}
}
{
}

\Verse{61}
{
    \zA{那婆子答应 去了。}
}
{
    \eA{Then pouring another cup, he was about to drain it,}
}
{
}

\Verse{62}
{
    \zA{不一时，只听得箫管悠扬，笙笛并发；正值风清气爽之时，那乐声穿林度水 而来，自然使人神怡心旷。}
}
{
    \eA{when he noticed that Madame Wang too was anxious for a drink, and that she
        bade a servant bring a warm supply of wine. "With alacrity, Pao-yü crossed
        over to her, and, presenting his own cup, he applied it to Madame Wang's
        lips. His mother drank two sips while he held it in his hands,}
}
{
}

\Verse{63}
{
    \zA{宝玉先禁不住，拿起壶来斟了一杯，一口饮尽，复又斟 上；才要饮，只见王夫人也要饮，命人换暖酒，宝玉连忙将自己的杯捧了过来，送
        到王夫人口边，王夫人便就他手内吃了两口。}
}
{
    \eA{but on the arrival of the warm wine, Pao-yü resumed his seat. Madame Wang
        laid hold of the warm decanter, and left the table, while the whole party
        quitted their places at the banquet; and Mrs. Hsüeh too rose to her feet.
        "Take over that decanter from her,"}
}
{
}

\Verse{64}
{
    \zA{一时暖酒来了，宝玉仍归旧坐。}
}
{
    \eA{dowager lady Chia promptly shouted to Li Wan and lady Feng, "and press your
        aunt into a seat. We shall all then feel at ease!"}
}
{
}

\Verse{65}
{
    \zA{王夫 人提了暖壶下席来，众人都出了席，薛姨妈也站起来，贾母忙命李凤二人接过壶来： “让你姨妈坐了，大家才便。”}
}
{
    \eA{Hearing this, Madame Wang surrendered the decanter to lady Feng and returned
        to her seat. "Let's all have a couple of cups of wine!" old lady Chia
        laughingly cried. "It's capital fun to-day!" With this proposal, she laid
        hold of a cup and offered it to Mrs. Hsüeh. Turning also towards Hsiang-yün
        and Pao-ch'ai:}
}
{
}

\Verse{66}
{
    \zA{王夫人见如此说，方将壶递与凤姐儿，自己归坐。}
}
{
    \eA{"You two cousins!" she added, "must also have a cup. Your cousin Lin can't
        take much wine, but even she mustn't be let off." While pressing them, she
        drained her cup.}
}
{
}

\Verse{67}
{
    \zA{贾母笑道：“大家吃上两杯，今日实在有趣。”}
}
{
    \eA{Hsiang-yün, Pao-ch'ai and Tai-y ü then had their drink. But about this time
        old goody Liu caught the strains of music, and, being already under the
        influence of liquor,}
}
{
}

\Verse{68}
{
    \zA{说着，擎杯让薛姨妈，又向湘云宝 钗道：“你姐妹两个也吃一杯。}
}
{
    \eA{her spirits became more and more exuberant, and she began to gesticulate and
        skip about. Her pranks amused Pao-yü to such a degree that leaving the
        table, he crossed over to where Tai-yü was seated and observed laughingly:}
}
{
}

\Verse{69}
{
    \zA{你林妹妹不大会吃，也别饶他。”}
}
{
    \eA{"Just you look at the way old goody Liu is going on!" "In days of yore,"
        Tai-yü smiled,}
}
{
}

\Verse{70}
{
    \zA{说着自己也干了， 湘云、宝钗、黛玉也都吃了。}
}
{
    \eA{"every species of animal commenced to dance the moment the sounds of music
        broke forth. She's like a buffalo now."}
}
{
}

\Verse{71}
{
    \zA{当下刘老老听见这般音乐，且又有了酒，越发喜的手 舞足蹈起来。}
}
{
    \eA{This simile made her cousins laugh. But shortly the music ceased. "We've all
        had our wine," Mrs. Hsüeh smilingly proposed, "so let's go and stroll about
        for a time;}
}
{
}

\Verse{72}
{
    \zA{宝玉因下席过来，向黛玉笑道：“你瞧刘老老的样子。”}
}
{
    \eA{we can after that sit down again!" Dowager lady Chia herself was at the
        moment feeling a strong inclination to have a ramble. In due course,
        therefore,}
}
{
}

\Verse{73}
{
    \zA{黛玉笑道： “当日圣乐一奏，百兽率舞，如今才一牛耳。”}
}
{
    \eA{they all left the banquet and went with their old senior, for a walk.
        Dowager lady Chia, however, longed to take goody Liu along with her to help
        her dispel her ennui,}
}
{
}

\Verse{74}
{
    \zA{众姐妹都笑了。}
}
{
    \eA{so promptly seizing the old dame's hand in hers,}
}
{
}

\Verse{75}
{
    \zA{须臾乐止，薛姨妈笑道：“大家的酒也都有了，且出去散散再坐罢。”}
}
{
    \eA{they threaded their way as far as the trees, which stood facing the hill.
        After lolling about with her for a few minutes, "What kind of tree is this?"
        she went on to inquire of her.}
}
{
}

\Verse{76}
{
    \zA{贾母也 正要散散，于是大家出席，都随着贾母游玩。}
}
{
    \eA{"What kind of stone is this? What species of flower is that?" Old goody Liu
        gave suitable reply to each of her questions. "Who'd ever have imagined it,"}
}
{
}

\Verse{77}
{
    \zA{贾母因要带着刘老老散闷，遂携了刘 老老至山前树下，盘桓了半晌，又说给他这是什么树，这是什么石，这是什么花。}
}
{
    \eA{she proceeded to tell dowager lady Chia; "not only are the human beings in
        the city grand, but even the birds are grand. Why, the moment these birds
        fly into your mansion, they also become beautiful things, and acquire the
        gift of speech as well!" The company could not make out the drift of her
        observations.}
}
{
}

\Verse{78}
{
    \zA{刘老老一一领会，又向贾母道：“谁知城里不但人尊贵，连雀儿也是尊贵的。}
}
{
    \eA{"What birds get transformed into beautiful things and become able to speak?"
        they felt impelled to ask. "Those perched on those gold stands, under the
        verandah, with green plumage and red beaks are parrots.}
}
{
}

\Verse{79}
{
    \zA{偏这 雀儿到了你们这里，他也变俊了，也会说话了。”}
}
{
    \eA{I know them well enough!" Goody Liu replied. "But those old black crows in
        the cages there have crests like phoenixes! They can talk too!" One and all
        laughed.}
}
{
}

\Verse{80}
{
    \zA{众人不解，因问：“什么雀儿变 俊了会说话？”}
}
{
    \eA{But not long elapsed before they caught sight of several waiting-maids, who
        came to invite them to a collation.}
}
{
}

\Verse{81}
{
    \zA{刘老老道：“那廊上金架子上站的绿毛红嘴是鹦哥儿，我是认得的。}
}
{
    \eA{"After the number of cups of wine I've had," old lady Chia said, "I don't
        feel hungry. But never mind, bring the things here. We can nibble something
        at our leisure."}
}
{
}

\Verse{82}
{
    \zA{那笼子里的黑老鸹子，又长出凤头儿来，也会说话呢！”}
}
{
    \eA{The maids speedily went off and fetched two teapoys; but they also brought a
        couple of small boxes with partitions. When they came to be opened and to be
        examined,}
}
{
}

\Verse{83}
{
    \zA{众人听了又都笑起来。}
}
{
    \eA{the contents of each were found to consist of two kinds of viands.}
}
{
}

\Verse{84}
{
    \zA{一时只见丫头们来请用点心，贾母道：“吃了两杯酒，倒也不饿。}
}
{
    \eA{In the one, were two sorts of steamed eatables. One of these was a sweet
        cake, made of lotus powder, scented with sun-flower. The other being rolls
        with goose fat and fir cone seeds.}
}
{
}

\Verse{85}
{
    \zA{也罢，就拿 了来这里，大家随便吃些罢。”}
}
{
    \eA{The second box contained two kinds of fried eatables; one of which was small
        dumplings,}
}
{
}

\Verse{86}
{
    \zA{丫头听说，便去抬了两张几来，又端了两个小捧盒。}
}
{
    \eA{about an inch in size. "What stuffing have they put in them?" dowager lady
        Chia asked. "They're with crabs inside," 'hastily rejoined the matrons.}
}
{
}

\Verse{87}
{
    \zA{揭开看时，每个盒内两样。}
}
{
    \eA{Their old mistress, at this reply, knitted her eyebrows. "These fat,}
}
{
}

\Verse{88}
{
    \zA{这盒内是两样蒸食：一样是藕粉桂花糖糕，一样是松瓤 鹅油卷。}
}
{
    \eA{greasy viands for such a time!" she observed. "Who'll ever eat these
        things?" But finding, when she came to inspect the other kind, that it
        consisted of small fruits of flour,}
}
{
}

\Verse{89}
{
    \zA{那盒内是两样炸的：一样是只有一寸来大的小饺儿。}
}
{
    \eA{fashioned in every shape, and fried in butter, she did not fancy these
        either. She then however pressed Mrs. Hsüeh to have something to eat,}
}
{
}

\Verse{90}
{
    \zA{贾母因问：“什么馅 子？”}
}
{
    \eA{but Mrs. Hsüeh merely took a piece of cake, while dowager lady Chia helped
        herself to a roll;}
}
{
}

\Verse{91}
{
    \zA{婆子们忙回：“是螃蟹的。”}
}
{
    \eA{but after tasting a bit, she gave the remaining half to a servant girl.}
}
{
}

\Verse{92}
{
    \zA{贾母听了，皱眉说道：“这会子油腻腻的，谁 吃这个。”}
}
{
    \eA{Goody Liu saw how beautifully worked those small flour fruits were, made as
        they were in various colours and designs, and she took, after picking and
        choosing,}
}
{
}

\Verse{93}
{
    \zA{又看那一样，是奶油炸的各色小面果子。}
}
{
    \eA{one which looked like a peony. "The most ingenious girls in our village
        could not, even with a pair of scissors,}
}
{
}

\Verse{94}
{
    \zA{也不喜欢，因让薛姨妈，薛姨 妈只拣了块糕。}
}
{
    \eA{cut out anything like this in paper!" she exclaimed. "I would like to eat
        it, but I can't make up my mind to!}
}
{
}

\Verse{95}
{
    \zA{贾母拣了个卷子，只尝了一尝，剩的半个，递给丫头了。}
}
{
    \eA{I had better pack up a few and take them home and give them to them as
        specimens!" Her remarks amused every one. "When you start for home,"}
}
{
}

\Verse{96}
{
    \zA{刘老老因 见那小面果子儿都玲珑剔透，各式各样，又拣了一朵牡丹花样的，笑道：“我们乡 里最巧的姐儿们，剪子也不能铰出这么个纸的来。}
}
{
    \eA{dowager lady Chia said, "I'll give you a whole porcelain jar full of them;
        so you may as well eat these first, while they are hot!" The rest of the
        inmates selected such of the fruits as took their fancy, but after they had
        helped themselves to one or two, they felt satisfied. Goody Liu, however,
        had never before touched such delicacies.}
}
{
}

\Verse{97}
{
    \zA{我又爱吃，又舍不得吃，包些家 去给他们做花样子去倒好。”}
}
{
    \eA{These were, in addition, made small, dainty, and without the least semblance
        of clumsiness, so when she and Pan Erh had served themselves to a few of
        each sort, half the contents of the dish vanished.}
}
{
}

\Verse{98}
{
    \zA{众人都笑了。}
}
{
    \eA{But what remained of them were then,}
}
{
}

\Verse{99}
{
    \zA{贾母笑道：“家去我送你一磁坛子，你 先趁热吃罢。”}
}
{
    \eA{at the instance of lady Feng, put into two plates, and sent, together with a
        partition-box, to Wen Kuan and the other singing girls as their share.}
}
{
}

\Verse{100}
{
    \zA{别人不过拣各人爱吃的拣了一两样就算了，刘老老原不曾吃过这些 东西，且都做的小巧，不显堆垛儿，他和板儿每样吃了些个，就去了半盘子。}
}
{
    \eA{At an unexpected moment, they perceived the nurse come in with Ta Chieh-erh
        in her arms, and they all induced her to have a romp with them for a time.
        But while Ta Chieh-erh was holding a large pumelo and amusing herself with
        it, she casually caught sight of Pan Erh with a 'Buddha's hand.' Ta Chieh
        would have it. A servant-girl endeavoured to coax (Pan-Erh) to surrender it
        to her,}
}
{
}

\Verse{101}
{
    \zA{剩的， 凤姐又命攒了两盘，并一个攒盒，给文官儿等吃去。}
}
{
    \eA{but Ta Chieh-erh, unable to curb her impatience, burst out crying. It was
        only after the pumelo had been given to Pan-Erh, and that the 'Buddha's
        hand' had, by dint of much humouring,}
}
{
}

\Verse{102}
{
    \zA{忽见奶子抱了大姐儿来，大家哄他玩了一会。}
}
{
    \eA{been got from Pan Erh and given to her, that she stopped crying. Pan Erh had
        played quite long enough with the 'Buddha's hand,'}
}
{
}

\Verse{103}
{
    \zA{那大姐儿因抱着一个大柚子玩， 忽见板儿抱着一个佛手，大姐儿便要。}
}
{
    \eA{and had, at the moment, his two hands laden with fruits, which he was in the
        course of eating. When he suddenly besides saw how scented and round the
        pumelo was,}
}
{
}

\Verse{104}
{
    \zA{丫鬟哄他取去，大姐儿等不得，便哭了。}
}
{
    \eA{the idea dawned on him that it was more handy for play, and, using it as a
        ball, he kicked it along and went off to have some fun,}
}
{
}

\Verse{105}
{
    \zA{众 人忙把柚子给了板儿，将板儿的佛手哄过来给他才罢。}
}
{
    \eA{relinquishing at once every thought of the 'Buddha's hand.' By this time
        dowager lady Chia and the other members had had tea, so leading off again
        goody Liu,}
}
{
}

\Verse{106}
{
    \zA{那板儿因玩了半日佛手，此 刻又两手抓着些果子吃，又见这个柚子又香又圆，更觉好玩，且当球踢着玩去，也 就不要佛手了。}
}
{
    \eA{they threaded their way to the Lung Ts'ui monastery. Miao Yü hastened to
        usher them in. On their arrival in the interior of the court, they saw the
        flowers and trees in luxuriant blossom. "Really," smiled old lady Chia,
        "it's those people, who devote themselves to an ascetic life and have
        nothing to do,}
}
{
}

\Verse{107}
{
    \zA{当下贾母等吃过了茶，又带了刘老老至栊翠庵来。}
}
{
    \eA{who manage, by constant repairs, to make their places much nicer than those
        of others!" As she spoke, she wended her steps towards the Eastern hall.}
}
{
}

\Verse{108}
{
    \zA{妙玉相迎进去。}
}
{
    \eA{Miao Yü, with a face beaming with smiles,}
}
{
}

\Verse{109}
{
    \zA{众人至院中， 见花木繁盛，贾母笑道：“到底是他们修行的人，没事常常修理，比别处越发好看。”}
}
{
    \eA{made way for her to walk in. "We've just been filling ourselves with wines
        and meats," dowager lady Chia observed, "and with the josses you've got in
        here, we shall be guilty of profanity. We'd better therefore sit here! But
        give us some of that good tea of yours;}
}
{
}

\Verse{110}
{
    \zA{一面说，一面便往东禅堂来。}
}
{
    \eA{and we'll get off so soon as we have had a cup of it." Pao-yü watched Miao
        Yü's movements intently,}
}
{
}

\Verse{111}
{
    \zA{妙玉笑往里让，贾母道：“我们才都吃了酒肉，你这 里头有菩萨，冲了罪过。}
}
{
    \eA{when he noticed her lay hold of a small tea-tray, fashioned in the shape of
        a peony, made of red carved lacquer, and inlaid with designs in gold
        representing a dragon ensconced in the clouds with the character 'longevity'
        clasped in its jaws,}
}
{
}

\Verse{112}
{
    \zA{我们这里坐坐，把你的好茶拿来，我们吃一杯就去了。”}
}
{
    \eA{a tray, which contained a small multicoloured cup with cover, fabricated at
        the 'Ch'eng' Kiln, and present it to his grandmother. "I don't care for 'Liu
        An' tea!"}
}
{
}

\Verse{113}
{
    \zA{宝玉留神看他是怎么行事，只见妙玉亲自捧了一个海棠花式雕漆填金“云龙献寿” 的小茶盘，里面放一个成窑五彩小盖钟，捧与贾母。}
}
{
    \eA{old lady Chia exclaimed. "I know it; but this is old 'Chün Mei,'" Miao Yü
        answered with a smile. Dowager lady Chia received the cup. "What water is
        this?" she went on to inquire.}
}
{
}

\Verse{114}
{
    \zA{贾母道：“我不吃六安茶。”}
}
{
    \eA{"It's rain water collected last year;" Miao Yü added by way of reply.}
}
{
}

\Verse{115}
{
    \zA{妙玉笑说：“知道。}
}
{
    \eA{Old lady Chia readily drank half a cup of the tea;}
}
{
}

\Verse{116}
{
    \zA{这是‘老君眉’。”}
}
{
    \eA{and smiling, she proffered it to goody Liu.}
}
{
}

\Verse{117}
{
    \zA{贾母接了，又问：“是什么水？”}
}
{
    \eA{"Just you taste this tea!" she said. Goody Liu drained the remainder with
        one draught.}
}
{
}

\Verse{118}
{
    \zA{妙玉道： “是旧年蠲的雨水。”}
}
{
    \eA{"It's good, of course," she remarked laughingly, "but it's rather weak! It
        would be far better were it brewed a little stronger!"}
}
{
}

\Verse{119}
{
    \zA{贾母便吃了半盏，笑着递与刘老老，说：“你尝尝这个茶。”}
}
{
    \eA{Dowager lady Chia and all the inmates laughed. But subsequently, each of
        them was handed a thin, pure white covered cup, all of the same make,
        originating from the 'Kuan' kiln.}
}
{
}

\Verse{120}
{
    \zA{刘老老便一口吃尽，笑道：“好是好，就是淡些，再熬浓些更好了。”}
}
{
    \eA{Miao Yü, however, soon gave a tug at Pao-ch'ai's and Tai-yü's lapels, and
        both quitted the apartment along with her. But Pao-yü too quietly followed
        at their heels.}
}
{
}

\Verse{121}
{
    \zA{贾母众人都 笑起来。}
}
{
    \eA{Spying Miao Yü show his two cousins into a side-room,}
}
{
}

\Verse{122}
{
    \zA{然后众人都是一色的官窑脱胎填白盖碗。}
}
{
    \eA{Pao-ch'ai take a seat in the court, Tai-yü seat herself on Miao Yü's rush
        mat, and Miao Yü herself approach a stove,}
}
{
}

\Verse{123}
{
    \zA{那妙玉便把宝钗黛玉的衣襟一拉，二人随他出去。}
}
{
    \eA{fan the fire and boil some water, with which she brewed another pot of tea,
        Pao-yü walked in. "Are you bent upon drinking your own private tea?"}
}
{
}

\Verse{124}
{
    \zA{宝玉悄悄的随后跟了来。}
}
{
    \eA{he smiled. "Here you rush again to steal our tea," the two girls laughed
        with one accord.}
}
{
}

\Verse{125}
{
    \zA{只 见妙玉让他二人在耳房内，宝钗便坐在榻上，黛玉便坐在妙玉的蒲团上。}
}
{
    \eA{"There's none for you!" But just as Miao Yü was going to fetch a cup, she
        perceived an old taoist matron bring away the tea things, which had been
        used in the upper rooms. "Don't put that 'Ch'eng' kiln tea-cup by!"}
}
{
}

\Verse{126}
{
    \zA{妙玉自向 风炉上煽滚了水，另泡了一壶茶。}
}
{
    \eA{Miao Yü hastily shouted. "Go and put it outside!" Pao-yü understood that it
        must be because old goody Liu had drunk out of it that she considered it too
        dirty to keep.}
}
{
}

\Verse{127}
{
    \zA{宝玉便轻轻走进来，笑道：“你们吃体己茶呢！”}
}
{
    \eA{He then saw Miao Yü produce two other cups. The one had an ear on the side.
        On the bowl itself were engraved in three characters:}
}
{
}

\Verse{128}
{
    \zA{二人都笑道：“你又赶了来撤茶吃!这里并没你吃的。”}
}
{
    \eA{'calabash cup,' in the plain 'square' writing. After these, followed a row
        of small characters in the 'true' style, to the effect that the cup had been
        an article much treasured by Wang K'ai.}
}
{
}

\Verse{129}
{
    \zA{妙玉刚要去取杯，只见道 婆收了上面茶盏来，妙玉忙命：“将那成窑的茶杯别收了，搁在外头去罢。”}
}
{
    \eA{Next came a second row of small characters stating: 'that in the course of
        the fourth moon of the fifth year of Yuan Feng, of the Sung dynasty, Su Shih
        of Mei Shan had seen it in the 'Secret' palace. This cup, Miao Yü filled,
        and handed to Pao-ch'ai. The other cup was, in appearance,}
}
{
}

\Verse{130}
{
    \zA{宝玉 会意，知为刘老老吃了，他嫌腌？}
}
{
    \eA{as clumsy as it was small; yet on it figured an engraved inscription,
        consisting of 'spotted rhinoceros cup,'}
}
{
}

\Verse{131}
{
    \zA{不要了。}
}
{
    \eA{in three 'seal' characters,}
}
{
}

\Verse{132}
{
    \zA{又见妙玉另拿出两只杯来，一个旁边有 一耳，杯上镌着“？？？”}
}
{
    \eA{which bore the semblance of pendent pearls. Miao Yü replenished this cup and
        gave it to Tai-yü; and taking the green jade cup, which she had, on previous
        occasions,}
}
{
}

\Verse{133}
{
    \zA{三个隶字，后有一行小真字，是“王恺珍玩”；又有“宋 元丰五年四月眉山苏轼见于秘府”一行小字。}
}
{
    \eA{often used for her own tea, she filled it and presented it to Pao-yü. "'The
        rules observed in the world,' the adage says, 'must be impartial,'" Pao-yü
        smiled. "But while my two cousins are handling those antique and rare gems,}
}
{
}

\Verse{134}
{
    \zA{妙玉斟了一？}
}
{
    \eA{here am I with this coarse object!"}
}
{
}

\Verse{135}
{
    \zA{递与宝钗。}
}
{
    \eA{"Is this a coarse thing?"}
}
{
}

\Verse{136}
{
    \zA{那一只形似 钵而小，也有三个垂珠篆字，镌着“点犀？”}
}
{
    \eA{Miao Yü exclaimed. "Why, I'm making no outrageous statement when I say that
        I'm inclined to think that it is by no means certain that you could lay your
        hand upon any such coarse thing as this in your home!"}
}
{
}

\Verse{137}
{
    \zA{。妙玉斟了一？}
}
{
    \eA{"'Do in the country as country people do,'}
}
{
}

\Verse{138}
{
    \zA{与黛玉，仍将前番自 己常日吃茶的那只绿玉斗来斟与宝玉。}
}
{
    \eA{the proverb says," Pao-yü laughingly rejoined. "So when one gets in a place
        like this of yours, one must naturally look down upon every thing in the way
        of gold,}
}
{
}

\Verse{139}
{
    \zA{宝玉笑道：“常言‘世法平等’：他两个就 用那样古玩奇珍，我就是个俗器了？”}
}
{
    \eA{pearls, jade and precious stones, as coarse rubbish!" This sentiment highly
        delighted Miao Yü. So much so, that producing another capacious cup, carved
        out of a whole bamboo root, which with its nine curves and ten rings,}
}
{
}

\Verse{140}
{
    \zA{妙玉道：“这是俗器？}
}
{
    \eA{with twenty knots in each ring, resembled a coiled dragon,}
}
{
}

\Verse{141}
{
    \zA{不是我说狂话，只怕 你家里未必找的出这么一个俗器来呢！”}
}
{
    \eA{"Here," she said with a face beaming with smiles, "there only remains this
        one! Can you manage this large cup?" "I can!" Pao-yü vehemently replied,}
}
{
}

\Verse{142}
{
    \zA{宝玉笑道：“俗语说：‘随乡入乡’，到 了你这里，自然把这金珠玉宝一概贬为俗器了。”}
}
{
    \eA{with high glee. "Albeit you have the stomach to tackle all it holds," Miao
        Yü laughed, "I haven't got so much tea for you to waste! Have you not heard
        how that the first cup is the 'taste'-cup; the second 'the stupid-thing-for-
        quenching-one's-thirst,'}
}
{
}

\Verse{143}
{
    \zA{妙玉听如此说，十分欢喜，遂又 寻出一只九曲十环一百二十节蟠虬整雕竹根的一个大盏出来，笑道：“就剩了这一 个，你可吃的了这一海？”}
}
{
    \eA{and the third 'the drink-mule' cup? But were you now to go in for this huge
        cup, why what more wouldn't that be?" At these words, Pao-ch'ai, Tai-yü and
        Pao-yü simultaneously indulged in laughter. But Miao-yü seized the teapot,
        and poured well-nigh a whole cupful of tea into the big cup.}
}
{
}

\Verse{144}
{
    \zA{宝玉喜的忙道：“吃的了。”}
}
{
    \eA{Pao-yü tasted some carefully, and found it, in real truth,}
}
{
}

\Verse{145}
{
    \zA{妙玉笑道：“你虽吃的了， 也没这些茶你遭塌。}
}
{
    \eA{so exceptionally soft and pure that he extolled it with incessant praise.
        "If you've had any tea this time," Miao-Yü pursued with a serious expression
        about her face,}
}
{
}

\Verse{146}
{
    \zA{岂不闻一杯为品，二杯即是解渴的蠢物，三杯便是饮驴了。}
}
{
    \eA{"it's thanks to these two young ladies; for had you come alone, I wouldn't
        have given you any." "I'm well aware of this,"}
}
{
}

\Verse{147}
{
    \zA{你 吃这一海，更成什么？”}
}
{
    \eA{Pao-yü laughingly rejoined, "so I too will receive no favour from your
        hands,}
}
{
}

\Verse{148}
{
    \zA{说的宝钗、黛玉、宝玉都笑了。}
}
{
    \eA{but simply express my thanks to these two cousins of mine,}
}
{
}

\Verse{149}
{
    \zA{妙玉执壶，只向海内斟了 约有一杯。}
}
{
    \eA{and have done!" "What you say makes your meaning clear enough!" Miao-yü
        said, when she heard his reply.}
}
{
}

\Verse{150}
{
    \zA{宝玉细细吃了，果觉轻淳无比，赏赞不绝。}
}
{
    \eA{"Is this rain water from last year?" Tai-yü then inquired. "How is it,"
        smiled Miao Yü sardonically,}
}
{
}

\Verse{151}
{
    \zA{妙玉正色道：“你这遭吃茶， 是托他两个的福，独你来了，我是不能给你吃的。”}
}
{
    \eA{"that a person like you can be such a boor as not to be able to discriminate
        water, when you taste it? This is snow collected from the plum blossom, five
        years back, when I was in the P'an Hsiang temple at Hsüan Mu.}
}
{
}

\Verse{152}
{
    \zA{宝玉笑道：“我深知道，我也 不领你的情，只谢他二人便了。”}
}
{
    \eA{All I got was that flower jar, green as the devil's face, full, and as I
        couldn't make up my mind to part with it and drink it, I interred it in the
        ground,}
}
{
}

\Verse{153}
{
    \zA{妙玉听了，方说：“这话明白。”}
}
{
    \eA{and only opened it this summer. I've had some of it once before, and this is
        the second time.}
}
{
}

\Verse{154}
{
    \zA{黛玉因问：“这也是旧年的雨水？”}
}
{
    \eA{But how is it you didn't detect it, when you put it to your lips? Has rain
        water, obtained a year back,}
}
{
}

\Verse{155}
{
    \zA{妙玉冷笑道：“你这么个人，竟是大俗人， 连水也尝不出来!这是五年前我在玄墓蟠香寺住着，收的梅花上的雪，统共得了那
        一鬼脸青的花瓮一瓮，总舍不得吃，埋在地下，今年夏天才开了。}
}
{
    \eA{ever got such a soft and pure flavour? and how possibly could it be drunk at
        all?" Tai-yü knew perfectly what a curious disposition she naturally had,
        and she did not think it advisable to start any lengthy discussion with her.
        Nor did she feel justified to protract her stay,}
}
{
}

\Verse{156}
{
    \zA{我只吃过一回， 这是第二回了。}
}
{
    \eA{so after sipping her tea, she intimated to Pao-ch'ai her intention to go,}
}
{
}

\Verse{157}
{
    \zA{你怎么尝不出来？}
}
{
    \eA{and they quitted the apartment. Pao-yü gave a forced smile to Miao Yü.}
}
{
}

\Verse{158}
{
    \zA{隔年蠲的雨水，那有这样清淳？}
}
{
    \eA{"That cup," he said, "is, of course, dirty; but is it not a pity to put it
        away for no valid reason?}
}
{
}

\Verse{159}
{
    \zA{如何吃得！”}
}
{
    \eA{To my idea it would be preferable,}
}
{
}

\Verse{160}
{
    \zA{宝钗 知他天性怪僻，不好多话，亦不好多坐，吃过茶，便约着黛玉走出来。}
}
{
    \eA{wouldn't it? to give it to that poor old woman; for were she to sell it, she
        could have the means of subsistence! What do you say, will it do?" Miao Yü
        listened to his suggestion, and then nodded her head,}
}
{
}

\Verse{161}
{
    \zA{宝玉和妙玉 陪笑说道：“那茶杯虽然腌？}
}
{
    \eA{after some reflection. "Yes, that will be all right!" she answered. "Lucky
        for her I've never drunk a drop out of that cup, for had I, I would rather
        have smashed it to atoms than have let her have it!}
}
{
}

\Verse{162}
{
    \zA{了，白撩了岂不可惜？}
}
{
    \eA{If you want to give it to her, I don't mind a bit about it;}
}
{
}

\Verse{163}
{
    \zA{依我说，不如就给了那贫婆子 罢，他卖了也可以度日。}
}
{
    \eA{but you yourself must hand it to her! Now, be quick and clear it away at
        once!" "Of course; quite so!" Pao-yü continued. "How could you ever go and
        speak to her?}
}
{
}

\Verse{164}
{
    \zA{你说使得么？”}
}
{
    \eA{Things would then come to a worse pass.}
}
{
}

\Verse{165}
{
    \zA{妙玉听了，想了一想，点头说道：“这也 罢了。}
}
{
    \eA{You too would be contaminated! If you give it to me, it will be all right."
        Miao Yü there and then directed some one to fetch it and to give it to
        Pao-yü.}
}
{
}

\Verse{166}
{
    \zA{幸而那杯子是我没吃过的；若是我吃过的，我就砸碎了也不能给他。}
}
{
    \eA{When it was brought, Pao-yü took charge of it. "Wait until we've gone out,"
        he proceeded, "and I'll call a few servant-boys and bid them carry several
        buckets of water from the stream and wash the floors;}
}
{
}

\Verse{167}
{
    \zA{你要给 他，我也不管，你只交给他快拿了去罢。”}
}
{
    \eA{eh, shall I?" "Yes, that would be better!" Miao Yü smiled. "The only thing
        is that you must tell them to bring the water,}
}
{
}

\Verse{168}
{
    \zA{宝玉道：“自然如此。}
}
{
    \eA{and place it outside the entrance door by the foot of the wall;}
}
{
}

\Verse{169}
{
    \zA{你那里和他说话 去？}
}
{
    \eA{for they mustn't come in." "This goes without saying!"}
}
{
}

\Verse{170}
{
    \zA{越发连你都腌？}
}
{
    \eA{Pao-yü said; and, while replying,}
}
{
}

\Verse{171}
{
    \zA{了。}
}
{
    \eA{he produced the cup from inside his sleeve,}
}
{
}

\Verse{172}
{
    \zA{只交给我就是了。”}
}
{
    \eA{and handed it to a young waiting-maid from dowager lady Chia's apartments to
        hold. "To-morrow," he told her,}
}
{
}

\Verse{173}
{
    \zA{妙玉便命人拿来递给宝玉。}
}
{
    \eA{"give this to goody Liu to take with her, when she starts on her way
        homewards!"}
}
{
}

\Verse{174}
{
    \zA{宝玉接了， 又道：“等我们出去了，我叫几个小么儿来河里打几桶水来洗地如何？”}
}
{
    \eA{By the time he made (the girl) understand the charge he entrusted her with,
        his old grandmother issued out and was anxious to return home. Miao Yü did
        not exert herself very much to induce her to prolong her visit;}
}
{
}

\Verse{175}
{
    \zA{妙玉笑道： “这更好了。}
}
{
    \eA{but seeing her as far the main gate, she turned round and bolted the doors.}
}
{
}

\Verse{176}
{
    \zA{只是你嘱咐他们，抬了水，只搁在山门外头墙根下，别进门来。”}
}
{
    \eA{But without devoting any further attention to her, we will now allude to
        dowager lady Chia. She felt thoroughly tired and exhausted. To such a
        degree, that she desired Madame Wang,}
}
{
}

\Verse{177}
{
    \zA{宝 玉道：“这是自然的。”}
}
{
    \eA{Ying Ch'un and her sisters to see that Mrs. Hsüeh had some wine,}
}
{
}

\Verse{178}
{
    \zA{说着，便袖着那杯递给贾母屋里的小丫头子拿着，说：“明 日刘老老家去，给他带去罢。”}
}
{
    \eA{while she herself retired to the Tao Hsiang village to rest. Lady Feng
        immediately bade some servants fetch a bamboo chair. On its arrival, dowager
        lady Chia seated herself in it, and two matrons carried her off hemmed in by
        lady Feng,}
}
{
}

\Verse{179}
{
    \zA{交代明白，贾母已经出来要回去。}
}
{
    \eA{Li Wan and a bevy of servant-girls, and matrons. But let us now leave her to
        herself,}
}
{
}

\Verse{180}
{
    \zA{妙玉亦不甚留， 送出山门，回身便将门闭了，不在话下。}
}
{
    \eA{without any additional explanations. During this while, Mrs. Hsüeh too said
        good bye and departed. Madame Wang then dismissed Wen Kuan and the other
        girls,}
}
{
}

\Verse{181}
{
    \zA{且说贾母因觉身上乏倦，便命王夫人和迎春姐妹陪着薛姨妈去吃酒，自己便往 稻香村来歇息。}
}
{
    \eA{and, distributing the eatables, that had been collected in the
        partition-boxes, to the servant-maids to go and feast on, she availed
        herself of the leisure moments to lie off; so reclining as she was,}
}
{
}

\Verse{182}
{
    \zA{凤姐忙命人将小竹椅抬来，贾母坐上，两个婆子抬起，凤姐李纨和 众丫头婆子围随去了，不在话下。}
}
{
    \eA{on the couch, which had been occupied by her old relative a few minutes
        back, she bade a young maid lower the portière; after which, she asked her
        to massage her legs. "Should our old lady yonder send any message, mind you
        call me at once,"}
}
{
}

\Verse{183}
{
    \zA{这里薛姨妈也就辞出。}
}
{
    \eA{she proceeded to impress on her mind, and, laying herself down,}
}
{
}

\Verse{184}
{
    \zA{王夫人打发文官等出去， 将攒盒散给众丫头们吃去，自己便也乘空歇着，随便歪在方才贾母坐的榻上，命一
        个小丫头放下帘子来，又命捶着腿，吩咐他：“老太太那里有信，你就叫我。”}
}
{
    \eA{she went to sleep. Pao-yü, Hsiang-yün and the rest watched the servant-girls
        take the partition-boxes and place them among the rocks, and seat themselves
        some on boulders, others on the turf-covered ground, some lean against the
        trees,}
}
{
}

\Verse{185}
{
    \zA{说 着也歪着睡着了。}
}
{
    \eA{others squat down besides the pool, and thoroughly enjoy themselves.}
}
{
}

\Verse{186}
{
    \zA{宝玉湘云等看着丫头们将攒盒搁在山石上，也有坐在山石上的， 也有坐在草地下的，也有靠着树的，也有傍着水的，倒也十分热闹。}
}
{
    \eA{But in a little time, they also perceived Yüan Yang arrive. Her object in
        coming was to carry off goody Liu for a stroll, so in a body they followed
        in their track, with a view of deriving some fun. Shortly, they got under
        the honorary gateway put up in the additional grounds, reserved for the
        imperial consort's visits to her parents,}
}
{
}

\Verse{187}
{
    \zA{一时又见鸳鸯来了，要带着刘老老逛，众人也都跟着取笑。}
}
{
    \eA{and old goody Liu shouted aloud: "Ai-yoh! What! Is there another big temple
        here!" While speaking, she prostrated herself and knocked her head, to the
        intense amusement of the company,}
}
{
}

\Verse{188}
{
    \zA{一时来至省亲别墅 的牌坊底下，刘老老道：“嗳呀!这里还有大庙呢。”}
}
{
    \eA{who were quite doubled up with laughter. "What are you laughing at?" goody
        Liu inquired. "I can decipher the characters on this honorary gateway. Over
        at our place temples of this kind are exceedingly plentiful;}
}
{
}

\Verse{189}
{
    \zA{说着，便爬下磕头。}
}
{
    \eA{and they've all got archways like this!}
}
{
}

\Verse{190}
{
    \zA{众人笑 弯了腰。}
}
{
    \eA{These characters give the name of the temple."}
}
{
}

\Verse{191}
{
    \zA{刘老老道：“笑什么？}
}
{
    \eA{"Can you make out from those characters what temple this is?"}
}
{
}

\Verse{192}
{
    \zA{这牌楼上的字我都认得。}
}
{
    \eA{they laughingly asked. Goody Liu quickly raised her head,}
}
{
}

\Verse{193}
{
    \zA{我们那里这样庙宇最多， 都是这样的牌坊，那字就是庙的名字。”}
}
{
    \eA{and, pointing at the inscription, "Are'nt these," she said, "the four
        characters 'Pearly Emperor's Precious Hall?'" Everybody laughed. They
        clapped their hands and applauded.}
}
{
}

\Verse{194}
{
    \zA{众人笑道：“你认得这是什么庙？”}
}
{
    \eA{But when about to chaff her again, goody Liu experienced a rumbling noise in
        her stomach,}
}
{
}

\Verse{195}
{
    \zA{刘老 老便抬头指那字道：“这不是‘玉皇宝殿’！”}
}
{
    \eA{and vehemently pulling a young servant-girl, and asking her for a couple of
        sheets of paper, she began immediately to loosen her garments.}
}
{
}

\Verse{196}
{
    \zA{众人笑的拍手打掌，还要拿他取笑 儿。}
}
{
    \eA{"It won't do in here!" one and all laughingly shouted out to her, and
        quickly they directed a matron to lead her away.}
}
{
}

\Verse{197}
{
    \zA{刘老老觉的肚里一阵乱响，忙的拉着一个丫头，要了两张纸，就解裙子。}
}
{
    \eA{When they got at the north-east corner, the matron pointed the proper place
        out to her, and in high spirits she walked off and went to have some rest.
        Goody Liu had taken plenty of wine;}
}
{
}

\Verse{198}
{
    \zA{众人 又是笑，又忙喝他：“这里使不得！”}
}
{
    \eA{she could not touch yellow wine; she had, what is more, drunk and eaten so
        many fat things that in the thirst,}
}
{
}

\Verse{199}
{
    \zA{忙命一个婆子，带了东北角上去了。}
}
{
    \eA{which supervened, she had emptied several cups of tea; the result was that
        she unavoidably got looseness of the bowels.}
}
{
}

\Verse{200}
{
    \zA{那婆子 指给他地方，便乐得走开去歇息。}
}
{
    \eA{She therefore squatted for ever so long before she felt any relief. But on
        her exit from the private chamber, the wind blew the wine to her head.}
}
{
}

\Verse{201}
{
    \zA{那刘老老因喝了些酒，他的脾气和黄酒不相宜， 且吃了许多油腻饮食发渴，多喝了几碗茶，不免通泻起来，蹲了半日方完。}
}
{
    \eA{Besides, being a woman well up in years, she felt, upon suddenly rising from
        a long squatting position, her eyes grow so dim and her head so giddy that
        she could not make out the way. She gazed on all four quarters, but the
        whole place being covered with trees, rockeries, towers, terraces,}
}
{
}

\Verse{202}
{
    \zA{及出厕 来，酒被风吹，且年迈之人，蹲了半天，忽一起身，只觉眼花头晕，辨不出路径。}
}
{
    \eA{and houses, she was quite at a loss how to determine her whereabouts, and
        where each road led to. She had no alternative but to follow a stone road,
        and to toddle on her way with leisurely step. But when she drew near a
        building,}
}
{
}

\Verse{203}
{
    \zA{四顾一望，都是树木山石，楼台房舍，却不知那一处是往那一路去的了，只得顺着 一条石子路慢慢的走来。}
}
{
    \eA{she could not make out where the door could be. After searching and
        searching, she accidentally caught sight of a bamboo fence. "Here's another
        trellis with flat bean plants creeping on it!" Goody Liu communed within
        herself. While giving way to reflection,}
}
{
}

\Verse{204}
{
    \zA{及至到了房子跟前又找不着门，再找了半日，忽见一带竹 篱。}
}
{
    \eA{she skirted the flower-laden hedge, and discovering a moonlike, cavelike,
        entrance, she stepped in. Here she discerned, stretching before her eyes a
        sheet of water,}
}
{
}

\Verse{205}
{
    \zA{刘老老心中自忖道：“这里也有扁豆架子？”}
}
{
    \eA{forming a pond, which measured no more than seven or eight feet in breadth.
        Its banks were paved with slabs of stone.}
}
{
}

\Verse{206}
{
    \zA{一面想，一面顺着花障走来，得 了个月洞门进去。}
}
{
    \eA{Its jadelike waves flowed in a limpid stream towards the opposite direction.
        At the upper end, figured a slab of white marble,}
}
{
}

\Verse{207}
{
    \zA{只见迎面一带水池，有七八尺宽，石头镶岸，里面碧波清水，上面有块白石横 架。}
}
{
    \eA{laid horizontally over the surface. Goody Liu wended her steps over the slab
        and followed the raised stone-road; then turning two bends, in the lake, an
        entrance into a house struck her gaze. Forthwith, she crossed the doorway,}
}
{
}

\Verse{208}
{
    \zA{刘老老便踱过石去，顺着石子甬路走去，转了两个弯子，只见有个房门。}
}
{
    \eA{but her eyes were soon attracted by a young girl, who advanced to greet her
        with a smile playing upon her lips. "The young ladies," goody Liu speedily
        remarked laughing, "have cast me adrift; they made me knock about,}
}
{
}

\Verse{209}
{
    \zA{于是 进了房门，便见迎面一个女孩儿，满面含笑的迎出来。}
}
{
    \eA{until I found my way in here." But seeing, after addressing her, that the
        girl said nothing by way of reply, goody Liu approached her and seized her
        by the hand,}
}
{
}

\Verse{210}
{
    \zA{刘老老忙笑道：“姑娘们把 我丢下了，叫我碰头碰到这里来了。”}
}
{
    \eA{when, with a crash, she fell against the wooden partition wall and bumped
        her head so that it felt quite sore. Upon close examination, she discovered
        that it was a picture.}
}
{
}

\Verse{211}
{
    \zA{说了，只觉那女孩儿不答。}
}
{
    \eA{"Do pictures really so bulge out!" Goody Liu mused within herself,}
}
{
}

\Verse{212}
{
    \zA{刘老老便赶来拉 他的手，咕咚一声却撞到板壁上，把头碰的生疼。}
}
{
    \eA{and, as she exercised her mind with these cogitations, she scanned it and
        rubbed her hand over it. It was perfectly even all over. She nodded her
        head, and heaved a couple of sighs.}
}
{
}

\Verse{213}
{
    \zA{细瞧了一瞧，原来是一幅画儿。}
}
{
    \eA{But the moment she turned round, she espied a small door over which hung a
        soft portière,}
}
{
}

\Verse{214}
{
    \zA{刘老老自忖道：“怎么画儿有这样凸出来的？”}
}
{
    \eA{of leek-green colour, bestrewn with embroidered flowers. Goody Liu lifted
        the portière and walked in. Upon raising her head, and casting a glance
        round,}
}
{
}

\Verse{215}
{
    \zA{一面想，一面看，一面又用手摸去， 却是一色平的，点头叹了两声。}
}
{
    \eA{she saw the walls, artistically carved in fretwork. On all four sides,
        lutes, double-edged swords, vases and censers were stuck everywhere over the
        walls;}
}
{
}

\Verse{216}
{
    \zA{一转身，方得了个小门，门上挂着葱绿撒花软帘， 刘老老掀帘进去。}
}
{
    \eA{and embroidered covers and gauze nets, glistened as brightly as gold, and
        shed a lustre vying with that of pearls. Even the bricks, on the ground, on
        which she trod, were jadelike green,}
}
{
}

\Verse{217}
{
    \zA{抬头一看，只见四面墙壁玲珑剔透，琴剑瓶炉皆贴在墙上，锦笼 纱罩，金彩珠光，连地下踩的砖皆是碧绿凿花，竟越发把眼花了。}
}
{
    \eA{inlaid with designs, so that her eyes got more and more dazzled. She tried
        to discover an exit, but where could she find a doorway? On the left, was a
        bookcase. On the right, a screen. As soon as she repaired behind the screen,
        she faced a door; but, she then caught sight of another old dame stepping in
        from outside,}
}
{
}

\Verse{218}
{
    \zA{找门出去，那里 有门？}
}
{
    \eA{and advancing towards her. Goody Liu was wonderstruck.}
}
{
}

\Verse{219}
{
    \zA{左一架书、右一架屏。}
}
{
    \eA{Her mind was full of uncertainty as to whether it might not be her
        son-in-law's mother.}
}
{
}

\Verse{220}
{
    \zA{刚从屏后得了一个门，只见一个老婆子也从外面迎着 进来。}
}
{
    \eA{"I expect," she felt prompted to ask with vehemence, "you went to the
        trouble of coming to hunt for me, as you didn't see me turn up at home for
        several days, eh? But what young lady introduced you in here?"}
}
{
}

\Verse{221}
{
    \zA{刘老老诧异，心中恍惚：莫非是他亲家母？}
}
{
    \eA{Then noticing that her whole head was bedecked with flowers, old goody Liu
        laughed. "How ignorant of the ways of the world you are!"}
}
{
}

\Verse{222}
{
    \zA{因问道：“你也来了，想是见我 这几日没家去？}
}
{
    \eA{she said. "Seeing the nice flowers in this garden, you at once set to work,
        forgetful of all consequences, and loaded your pate with them!"}
}
{
}

\Verse{223}
{
    \zA{亏你找我来，那位姑娘带进来的？”}
}
{
    \eA{However, while she derided her, the other old dame simply laughed, without
        making any rejoinder.}
}
{
}

\Verse{224}
{
    \zA{又见他戴着满头花，便笑道： “你好没见世面!见这里的花好，你就没死活戴了一头。”}
}
{
    \eA{But the recollection suddenly flashed to her memory that she had often heard
        of some kind of cheval-glasses, found in wealthy and well-to-do families,
        and, "May it not be," (she wondered), "my own self reflected in this glass!"}
}
{
}

\Verse{225}
{
    \zA{说着，那老婆子只是笑， 也不答言。}
}
{
    \eA{After concluding this train of thoughts, she put out her hands, and feeling
        it and then minutely scrutinising it,}
}
{
}

\Verse{226}
{
    \zA{刘老老便伸手去羞他的脸，他也拿手来挡，两个对闹着。}
}
{
    \eA{she realised that the four wooden partition walls were made of carved
        blackwood, into which mirrors had been inserted. "These have so far impeded
        my progress,"}
}
{
}

\Verse{227}
{
    \zA{刘老老一下子 却摸着了，但觉那老婆子的脸冰凉挺硬的，倒把刘老老唬了一跳。}
}
{
    \eA{she consequently exclaimed, "and how am I to manage to get out?" As she
        soliloquised, she kept on rubbing the mirror. This mirror was, in fact,
        provided with some western mechanism,}
}
{
}

\Verse{228}
{
    \zA{猛想起：“常听 见富贵人家有种穿衣镜，这别是我在镜子里头吗？”}
}
{
    \eA{which enabled it to open and shut, so while goody Liu inadvertently passed
        her hands, quite at random over its surface, the pressure happily fell on
        the right spot,}
}
{
}

\Verse{229}
{
    \zA{想毕，又伸手一抹，再细一看， 可不是四面雕空的板壁，将这镜子嵌在中间的，不觉也笑了。}
}
{
    \eA{and opening the contrivance, the mirror flung round, exposing a door to
        view. Old goody Liu was full of amazement as well as of admiration. With
        hasty step, she egressed. Her eyes unexpectedly fell on a most handsome set
        of bed-curtains.}
}
{
}

\Verse{230}
{
    \zA{因说：“这可怎么出 去呢？”}
}
{
    \eA{But being at the time still seven or eight tenths in the wind, and quite
        tired out from her tramp,}
}
{
}

\Verse{231}
{
    \zA{一面用手摸时，只听“硌磴”一声，又吓的不住的展眼儿。}
}
{
    \eA{she with one jump squatted down on the bed, saying to herself: "I'll just
        have a little rest." So little, however, did she, contrary to her
        expectations,}
}
{
}

\Verse{232}
{
    \zA{原来是西洋机 括，可以开合，不意刘老老乱摸之间，其力巧合，便撞开消息，掩过镜子，露出门 来。}
}
{
    \eA{have any control over herself, that, as she reeled backwards and forwards,
        her eyes got quite drowsy, and then the moment she threw herself in a
        recumbent position,}
}
{
}

\Verse{233}
{
    \zA{刘老老又惊又喜，遂走出来，忽见有一副最精致的床帐。}
}
{
    \eA{she dropped into a sound sleep. But let us now see what the others were up
        to. They waited for her and waited; but they saw nothing of her.}
}
{
}

\Verse{234}
{
    \zA{他此时又带了七八分 酒，又走乏了，便一屁股坐在床上。}
}
{
    \eA{Pan Erh got, in the absence of his grandmother, so distressed that he melted
        into tears. "May she not have fallen into the place?"}
}
{
}

\Verse{235}
{
    \zA{只说歇歇，不承望身不由己，前仰后合的，朦 胧两眼，一歪身就睡倒在床上。}
}
{
    \eA{one and all laughingly observed. "Be quick and tell some one to go and have
        a look!" Two matrons were directed to go in search of her; but they returned
        and reported that she was not to be found.}
}
{
}

\Verse{236}
{
    \zA{且说众人等他不见，板儿没了他老老，急的哭了。}
}
{
    \eA{The whole party instituted a search in every nook and corner, but nothing
        could be seen of her. "She was so drunk," Hsi Jen suggested,}
}
{
}

\Verse{237}
{
    \zA{众人都笑道：“别是掉在茅 厕里了？}
}
{
    \eA{"that she's sure to have lost her way, and following this road, got into our
        back-rooms.}
}
{
}

\Verse{238}
{
    \zA{快叫人去瞧瞧。”}
}
{
    \eA{Should she have crossed to the inner side of the hedge,}
}
{
}

\Verse{239}
{
    \zA{因命两个婆子去找。}
}
{
    \eA{she must have come to the door of the backhouse and got in.}
}
{
}

\Verse{240}
{
    \zA{回来说：“没有。”}
}
{
    \eA{Nevertheless, the young maids, she must have come across,}
}
{
}

\Verse{241}
{
    \zA{众人纳闷。}
}
{
    \eA{must know something about her.}
}
{
}

\Verse{242}
{
    \zA{还 是袭人想道：“一定他醉了，迷了路，顺着这条路往我们后院子里去了。}
}
{
    \eA{If she did not get inside the hedge, but continued in a south westerly
        direction, she's all right, if she made a detour and walked out. But if she
        hasn't done so,}
}
{
}

\Verse{243}
{
    \zA{要进了花 障子，打后门进去，还有小丫头子们知道；若不进花障子，再往西南上去，可够他 绕会子好的了!我瞧瞧去。”}
}
{
    \eA{why, she'll have enough of roaming for a good long while! I had better
        therefore go and see what she's up to." With these words still on her lips,
        she retraced her footsteps and repaired into the I Hung court.}
}
{
}

\Verse{244}
{
    \zA{说着便回来。}
}
{
    \eA{She called out to the servants,}
}
{
}

\Verse{245}
{
    \zA{进了怡红院，叫人，谁知那几个小丫头 已偷空玩去了。}
}
{
    \eA{but, who would have thought it, the whole bevy of young maids, attached to
        those rooms, had seized the opportunity to go and have a romp,}
}
{
}

\Verse{246}
{
    \zA{袭人进了房门，转过集锦？}
}
{
    \eA{so Hsi Jen straightway entered the door of the house. As soon as she turned
        the multicoloured embroidered screen,}
}
{
}

\Verse{247}
{
    \zA{子，就听的鼾？}
}
{
    \eA{the sound of snoring as loud as peals of thunder,}
}
{
}

\Verse{248}
{
    \zA{如雷，忙进来，只闻见酒屁臭气满 屋。}
}
{
    \eA{fell on her ear. Hastily she betook herself inside, but her nostrils were
        overpowered by the foul air of wine and w..d,}
}
{
}

\Verse{249}
{
    \zA{一瞧，只见刘老老扎手舞脚的仰卧在床上。}
}
{
    \eA{which infected the apartment. At a glance, she discovered old goody Liu
        lying on the bed, face downwards, with hands sprawled out and feet knocking
        about all over the place.}
}
{
}

\Verse{250}
{
    \zA{袭人这一惊不小，忙上来将他没死 活的推醒。}
}
{
    \eA{Hsi Jen sustained no small shock. With precipitate hurry, she rushed up to
        her, and, laying hold of her,}
}
{
}

\Verse{251}
{
    \zA{那刘老老惊醒，睁眼看见袭人，连忙爬起来，道：“姑娘，我该死了! 好歹并没弄腌？}
}
{
    \eA{lying as she was more dead than alive, she pushed her about until she
        succeeded in rousing her to her senses. Old goody Liu was startled out of
        her sleep. She opened wide her eyes, and, realising that Hsi Jen stood
        before her,}
}
{
}

\Verse{252}
{
    \zA{了床。”}
}
{
    \eA{she speedily crawled up.}
}
{
}

\Verse{253}
{
    \zA{一面说，用手去掸。}
}
{
    \eA{"Miss!" she pleaded. "I do deserve death! I have done what I shouldn't;}
}
{
}

\Verse{254}
{
    \zA{袭人恐惊动了宝玉，只向他摇手儿， 不叫他说话。}
}
{
    \eA{but I haven't in any way soiled the bed." So saying, she swept her hands
        over it. But Hsi Jen was in fear and trembling lest the suspicions of any
        inmate should be aroused,}
}
{
}

\Verse{255}
{
    \zA{忙将当地大鼎内贮了三四把百合香，仍用罩子罩上。}
}
{
    \eA{and lest Pao-yü should come to know of it, so all she did was to wave her
        hand towards her, bidding her not utter a word.}
}
{
}

\Verse{256}
{
    \zA{所喜不曾呕吐。}
}
{
    \eA{Then with alacrity grasping three or four handfuls of 'Pai Ho' incense,}
}
{
}

\Verse{257}
{
    \zA{忙悄悄的笑道：“不相干，有我呢。}
}
{
    \eA{she heaped it on the large tripod, which stood in the centre of the room,
        and put the lid back again;}
}
{
}

\Verse{258}
{
    \zA{你跟我出来罢。”}
}
{
    \eA{delighted at the idea that she had not been so upset as to be sick.}
}
{
}

\Verse{259}
{
    \zA{刘老老答应着，跟了袭人， 出至小丫头子们房中，命他坐下，因教他说道：“你说‘醉倒在山子石上，打了个 盹儿’就完了。”}
}
{
    \eA{"It doesn't matter!" she quickly rejoined in a low tone of voice with a
        smile, "I'm here to answer for this. Come along with me!" While old goody
        Liu expressed her readiness to comply with her wishes, she followed Hsi Jen
        out into the quarters occupied by the young maids.}
}
{
}

\Verse{260}
{
    \zA{刘老老答应“是”。}
}
{
    \eA{Here (Hsi Jen) desired her to take a seat.}
}
{
}

\Verse{261}
{
    \zA{又给了他两碗茶吃，方觉酒醒了。}
}
{
    \eA{"Mind you say," she enjoined her, "that you were so drunk that you stretched
        on a boulder and had a snooze!"}
}
{
}

\Verse{262}
{
    \zA{因问道： “这是那个小姐的绣房？}
}
{
    \eA{"All right! I will!" old goody Liu promised. Hsi Jen afterwards helped her
        to two cups of tea,}
}
{
}

\Verse{263}
{
    \zA{这么精致!我就像到了天宫里的似的。”}
}
{
    \eA{when she, at length, got over the effects of the wine. "What young lady's
        room is this that it is so beautiful?"}
}
{
}

\Verse{264}
{
    \zA{袭人微微的笑道： “这个么，是宝二爷的卧房啊。”}
}
{
    \eA{she then inquired. "It seemed to me just as if I had gone to the very
        heavenly palace." Hsi Jen gave a faint smile. "This one?" she asked. "Why,}
}
{
}

\Verse{265}
{
    \zA{那刘老老吓的不敢做声。}
}
{
    \eA{it's our master Secundus', Mr. Pao's bedroom." Old goody Liu was quite taken
        aback,}
}
{
}

\Verse{266}
{
    \zA{袭人带他从前面出去， 见了众人，只说：“他在草地下睡着了，带了他来的。”}
}
{
    \eA{and could not even presume to utter a sound. But Hsi Jen led her out across
        the front compound; and, when they met the inmates of the family, she simply
        explained to them that she had found her fast asleep on the grass,}
}
{
}

\Verse{267}
{
    \zA{众人都不理会，也就罢了。}
}
{
    \eA{and brought her along. No one paid any heed to the excuse she gave, and the
        subject was dropped.}
}
{
}

\Verse{268}
{
    \zA{一时贾母醒了，就在稻香村摆晚饭。}
}
{
    \eA{Presently, dowager lady Chia awoke, and the evening meal was at once served
        in the Tao Hsiang Ts'un.}
}
{
}

\Verse{269}
{
    \zA{贾母因觉懒懒的，也没吃饭，便坐了竹椅 小敞轿，回至房中歇息，命凤姐儿等去吃饭。}
}
{
    \eA{Dowager lady Chia was however quite listless, and felt so little inclined to
        eat anything that she forthwith got into a small open chair, with bamboo
        seat, and returned to her suite of rooms to rest.}
}
{
}

\Verse{270}
{
    \zA{他姐妹方复进园来。}
}
{
    \eA{But she insisted that lady Feng and her companions should go and have their
        repast,}
}
{
}

\Verse{271}
{
    \zA{未知如何，且看下回分解。}
}
{
    \eA{so the young ladies eventually adjourned once more into the garden. But,
        reader, you do not know the sequel,}
}
{
}

\Verse{272}
{
    \zA{(本章完)}
}
{
    \eA{so peruse the circumstances given in detail in the next chapter.}
}
{
}
